% TOP_PROBLEM: {"id": 20001067, "title": "Cut covers and small cases of Mader's local-connectivity conjecture", "category_id": 2, "category": {"id": 2, "name": "combinatorics", "display_name": "Combinatorics", "description": "Counting problems, graph theory, discrete structures.", "slug": "combinatorics", "order_index": 2, "created_at": "2026-07-31T15:26:25.671Z"}, "status": "partially_solved", "problem_number": "AIM-COMBINATORICS-0192", "set_id": 14, "statusQualification": "Loops are excluded to avoid the immediate one-vertex obstruction; parallel arcs retain the source general digraph convention."}
% FIRST_POSTED: 2026-10-01
% ENTRY_KIND: statement_only
% UnsolvedMath ID 20001067; source catalogue code AIM-COMBINATORICS-0192
% !TeX program = lualatex
% Definition and sources only; this file is not an LLM solution attempt.
\documentclass[11pt,a4paper]{article}
\usepackage[margin=25mm]{geometry}
\usepackage{fontspec}
\setmainfont{lmroman10-regular.otf}[BoldFont=lmroman10-bold.otf,
 ItalicFont=lmroman10-italic.otf,BoldItalicFont=lmroman10-bolditalic.otf]
\usepackage{amsmath,amssymb,mathtools}
\usepackage{unicode-math}
\setmathfont{latinmodern-math.otf}
\AtBeginDocument{\let\setminus\smallsetminus}
\usepackage{xurl,hyperref}
\hypersetup{colorlinks=true,urlcolor=blue}
\setcounter{secnumdepth}{0}
\setlength{\parindent}{0pt}
\setlength{\parskip}{5pt}
\setlength{\emergencystretch}{3em}
\begin{document}
\section*{Cut covers and small cases of Mader's local-connectivity conjecture}\label{20001067}
{\small UnsolvedMath ID 20001067 · AIM-COMBINATORICS-0192 · Combinatorics · AIM Workshop Problem Lists}
\subsection{Definitions and mathematical statement}
Let $D$ be a finite nonempty loopless digraph; parallel arcs are
allowed and counted separately. The outdegree $d^+(v)$ counts all
arcs with tail $v$, and $\delta^+(D)$ is their minimum over vertices.
For distinct vertices $x,y$, define $\lambda_D(x,y)$ to be the
largest number of pairwise edge-disjoint directed paths from $x$
to $y$. Paths may share internal vertices, but cannot use the same
arc; parallel arcs are distinct edges for this purpose. Each path
is directed consistently from $x$ to $y$.

The source's two assertions ask whether, for every integer $r\ge1$,
\[
 \delta^+(D)\ge r\quad\Longrightarrow\quad
 \exists x\ne y:\lambda_D(x,y)\ge r,
\]
and, more strongly, whether $x,y$ can always be chosen so that
$x\to y$ is an arc of $D$. The stronger assertion includes the first.
If an arc $x\to y$ is present, the one-edge path can be one member
of the family, but all $r$ paths must still be mutually edge-disjoint.
There is no prescribed pair of endpoints and no minimum-indegree
condition.
The direction in the conclusion is from the tail of the selected arc
to its head, rather than in the reverse direction. This is edge
connectivity, as denoted by $\lambda$ in the AIM source; replacing it
by internally vertex-disjoint connectivity would give a different
problem. Loops are excluded since they can raise outdegree without
providing a path to a different vertex.
\subsection{Short English statement}
Does minimum outdegree r force r edge-disjoint directed paths between some pair of vertices, and even between the ends of an arc?
\subsection{Sources}
\begin{itemize}
\item[S1] ULAM.AI and the UnsolvedMath Contributors, supplied problems.json, record 20001067 (AIM-COMBINATORICS-0192), AIM Workshop Problem Lists. Catalogue identity and source transcription; CC BY 4.0.
\url{https://huggingface.co/datasets/ulamai/UnsolvedMath}.
\item[S2] American Institute of Mathematics, The Caccetta-Haggkvist conjecture, item 6. Original workshop question underlying AIM-COMBINATORICS-0192.
\url{https://aimath.org/WWN/caccetta/caccetta.pdf}.
\end{itemize}
\end{document}
